
The Universe offers a unique precision laboratory for fundamental physics, carrying in its observables the imprints of processes at energy scales and in regimes that are inaccessible to any terrestrial experiment. This dissertation develops theoretical predictions and observational tools together, with the goal of turning current and upcoming measurements of the cosmic microwave background (CMB) and the large-scale structure (LSS) of the Universe into robust probes of fundamental physics. In particular, it addresses three interconnected frontiers of beyond-the-Standard-Model physics: the microphysics of inflation, the nature of the dark sector, and the free-streaming properties of cosmic neutrinos.

Reaching furthest back in cosmic history, inflation is usually assumed to take place in an empty, supercooled Universe, but a well-motivated alternative — warm inflation — posits that the inflaton remains thermally coupled to a bath of radiation throughout the accelerated expansion. We develop the theoretical framework and computational infrastructure needed to confront this scenario with data, and show that the persistent thermal environment gives rise to novel phenomena with no analogue in the standard picture, including a new mechanism for producing the observed dark matter abundance during inflation itself and an irreducible background of gravitational waves of thermal origin.

Closer to the present day, dark matter and dark energy together dominate the energy budget of the Universe, yet their fundamental nature remains unknown. We show that cosmological observations — in particular the CMB — and terrestrial experiments provide complementary and sometimes surprisingly powerful handles on dark matter properties, probing interaction strengths and candidate scenarios that neither approach could constrain alone. On the dark energy side, we demonstrate that the interpretation of recent observational hints for a time-varying equation of state depends critically on the theoretical assumptions underlying the analysis, and that incorporating well-motivated physical priors substantially weakens the apparent evidence for new physics.

Among the most elusive species in the Standard Model itself, cosmic neutrinos leave a wealth of imprints on cosmological observables, including a characteristic and coherent phase shift in the acoustic oscillations of the primordial plasma — a signature that is uniquely sensitive to their free-streaming nature and difficult to mimic through changes in other cosmological parameters. We develop a comprehensive framework for extracting this imprint from cosmological data, obtain a high-significance detection consistent with three free-streaming neutrino species, and extend the analysis to directly constrain non-standard neutrino interactions. We further identify a qualitatively new manifestation of this signature in the 21-cm signal emitted by neutral hydrogen during cosmic dawn, opening a complementary observational window with distinct sensitivity to the properties of neutrinos and other light relics.
